% RQ3 Background: Proposal Pipelines

\subsection{Proposal Pipelines in Practice}

\subsubsection{Compound}

Compound's Governor Bravo defines a linear pipeline:
\begin{enumerate}
    \item Proposal creation (requires token threshold)
    \item Voting period (3 days)
    \item Timelock (2 days)
    \item Execution
\end{enumerate}

No formal temp-check or fast-track; all proposals follow the same path.

\subsubsection{Uniswap}

Uniswap adds informal stages:
\begin{enumerate}
    \item Temperature check (Snapshot, 2 days)
    \item Consensus check (Snapshot, 5 days)
    \item Governance proposal (on-chain)
\end{enumerate}

This multi-stage approach filters proposals before expensive on-chain voting.

\subsubsection{Optimism}

Optimism's Token House uses:
\begin{enumerate}
    \item Draft proposal (feedback period)
    \item Review period
    \item Voting period
\end{enumerate}

Different proposal types have different requirements and timelines.

\subsubsection{Nouns DAO}

Nouns uses a streamlined single-stage process but with high proposal threshold (2 Nouns), effectively filtering at submission rather than through pipeline stages.

\subsection{Pipeline Design Considerations}

\subsubsection{Throughput}

Governance throughput measures proposals processed per unit time:
\begin{equation}
\text{throughput} = \frac{\text{proposals resolved}}{\text{time period}}
\end{equation}

High throughput enables responsive governance but may overwhelm participant attention.

\subsubsection{Decision Quality}

Quality is harder to define but includes:
\begin{itemize}
    \item Adequate deliberation time
    \item Sufficient participation
    \item Alignment with community preferences
    \item Implementation success rate
\end{itemize}

We use abandonment rate (proposals failing without resolution) as a quality proxy: high abandonment suggests proposals entering the pipeline that should have been filtered earlier.

\subsubsection{Time-to-Decision}

The interval from proposal creation to resolution:
\begin{equation}
\text{TTD}(p) = t_{\text{resolved}}(p) - t_{\text{created}}(p)
\end{equation}

Fast TTD enables agility; slow TTD may cause missed opportunities or stale decisions.

\subsection{Theoretical Foundations}

\subsubsection{Queuing Theory}

Proposal pipelines can be modeled as queuing systems. Arrival rate $\lambda$ (proposals per day) must be balanced against service rate $\mu$ (processing capacity):

\begin{equation}
\rho = \frac{\lambda}{\mu} < 1 \text{ for stable queue}
\end{equation}

When $\rho > 1$, proposal backlog grows unboundedly.

\subsubsection{Signal Filtering}

Temp-checks function as noisy classifiers, filtering proposals based on initial signal:
\begin{itemize}
    \item True positives: Good proposals that pass
    \item False positives: Poor proposals that pass (waste voting resources)
    \item True negatives: Poor proposals filtered
    \item False negatives: Good proposals incorrectly filtered
\end{itemize}

Threshold selection trades off false positives against false negatives.

\subsection{Prior Work}

Limited academic literature addresses DAO proposal pipelines specifically. \citet{snapshot2021governance} provides analytics on off-chain voting patterns. Our simulation framework enables systematic study of pipeline parameter effects not available from observational data alone.
